\scanpage{70}
$i=0$ seulmt) $C^i\xrightarrow{d^i}C^{i+1}$ ait un Ker \deleted{qui soit tel que $Z^i\longrightarrow C^i$ soit un $\Sigma$-mono (Ici, le $\Sigma$ dans $B$ est de la \unsure{poudre} aux yeux, il n'y a \unsure{à y penser}).} Ceci dit, on se borne à de tels complexes, qui correspondent donc à des foncteurs \deleted{\unsure{$\Sigma$-$P$}} $k$-add. \inserted{($\Sigma$-ex. à g)} $P_0\longrightarrow B$ resp. $Q\longrightarrow B$. \deleted{On se propose de \inserted{\unsure{se}} donner une structure multiplicative sur ceux-ci, ce qui nous oblige à donner de plus
\[
\deleted{F(\Psi^0)\otimes_B F(\mathbb{C}^i)\longrightarrow F(\Psi^0*\mathbb{C}^i)\text{ et }\unsure{\mathrm{sym}}}
\]
hom de bifoncteurs
\[
F(X)\otimes_B F(Y)\longrightarrow F(X*Y)
\]
sur $P_0\times P_0$ resp $Q\times Q$. Or \unsure{$P$}}

Si \unsure{enfin|même|tout} dans $B$ tte flèche a un noyau, alors on se bornera aux foncteurs $\Sigma$-ex. à g de $P$ dans $B$, en \deleted{\unsure{prolongeant}} interprétant ceux-ci comme des foncteurs $k$-additifs quelconques sur $\mathbb{K}$ i.e. des $C^{\bullet}\in\mathrm{ob}\,B^{\bullet}$. On peut l'expliciter ainsi: tt $L\in P$ \inserted{($\in Q,\in P_0$} \deleted{\unsure{(ou soit}} s'insère dans \deleted{deux} $\Sigma$-suites de $P$ ($Q,P_0$)
\[
\begin{aligned}
0\longrightarrow L\longrightarrow \unsure{I|\mathbb I}'\longrightarrow L'\longrightarrow0&\\
&0\longrightarrow L'\longrightarrow J^{\unsure{''|'}}\longrightarrow L''\longrightarrow0
\end{aligned}
\]
avec $I,J'\in\mathbb{K}$, et par suite en
% archive sheet number: 69
